\section{Control Flow in the Go Programming Language}

Control flow refers to the order in which statements are executed
during program execution. In the Go programming language, control
flow constructs determine how a program makes decisions, repeats
operations, selects among alternatives, and transfers execution
between different parts of a program. Go provides structured
control flow mechanisms, including conditional statements, loops,
switch statements, deferred function calls, and control transfer
statements.

These constructs enable developers to implement algorithms,
manage program state, handle exceptional run-time situations, and
express complex computational logic in a clear and maintainable
manner. Go does not provide a conventional \texttt{while} loop;
instead, iteration is implemented using the versatile
\texttt{for} statement.

\subsection{Conditional Statements}

Conditional statements allow a program to select which statements
to execute according to the evaluation of a Boolean expression.
Go provides the \texttt{if}, \texttt{else if}, and \texttt{else}
constructs for conditional execution. The condition specified
in an \texttt{if} statement must evaluate to a Boolean value.

The general syntax is as follows:

\begin{lstlisting}
if condition {
    // Statements executed when condition is true
} else if anotherCondition {
    // Statements executed when anotherCondition is true
} else {
    // Statements executed when all preceding conditions are false
}
\end{lstlisting}

The conditions are evaluated sequentially. When a condition
evaluates to \texttt{true}, its associated block is executed,
and the remaining branches are skipped. The \texttt{else} branch
is optional and provides an alternative execution path when no
preceding condition is satisfied.

For example:

\begin{lstlisting}
package main

import "fmt"

func main() {
    score := 75

    if score >= 80 {
        fmt.Println("Distinction")
    } else if score >= 60 {
        fmt.Println("Pass")
    } else {
        fmt.Println("Fail")
    }
}
\end{lstlisting}

In this example, the value of \texttt{score} determines which
message is displayed. Since the score is 75, the second
conditional branch is executed.

Go also supports an initialisation statement within an
\texttt{if} statement. A variable declared in this statement
is scoped to the conditional statement, including its
associated \texttt{else} branches.

\begin{lstlisting}
if value := calculate(); value > 0 {
    fmt.Println("Positive value")
} else {
    fmt.Println("Zero or negative value")
}
\end{lstlisting}

This feature allows temporary variables to be declared close
to the conditions in which they are used, thereby improving
encapsulation and reducing unnecessary scope.

\subsection{The \texttt{switch} Statement}

The \texttt{switch} statement provides a structured mechanism
for selecting one execution path from several alternatives.
It can be used when a variable or expression must be compared
against multiple possible values. Unlike a long sequence of
conditional statements, a \texttt{switch} statement can make
multi-way selection easier to read and maintain.

A general expression switch has the following form:

\begin{lstlisting}
switch expression {
case value1:
    // Statements for value1
case value2:
    // Statements for value2
default:
    // Statements when no case matches
}
\end{lstlisting}

Go evaluates the switch expression and compares its value
against the expressions in the \texttt{case} clauses. The
statements associated with the first matching case are
executed. The \texttt{default} clause is optional and is
executed when no case matches.

For example:

\begin{lstlisting}
package main

import "fmt"

func main() {
    day := 3

    switch day {
    case 1:
        fmt.Println("Monday")
    case 2:
        fmt.Println("Tuesday")
    case 3:
        fmt.Println("Wednesday")
    default:
        fmt.Println("Unknown day")
    }
}
\end{lstlisting}

The output is \texttt{Wednesday}, because the value of
\texttt{day} matches the third case.

Go does not automatically fall through from one case to
the next. Each case normally terminates after its statements
have executed. Explicit fall-through behaviour can be
requested using the \texttt{fallthrough} statement, which
transfers execution to the next case's statements without
evaluating that case's expressions.

A switch statement can also omit its expression. In this
form, each case is evaluated as a Boolean condition, making
the construct similar to a sequence of conditional branches.

\begin{lstlisting}
temperature := 28

switch {
case temperature < 0:
    fmt.Println("Freezing conditions")
case temperature < 20:
    fmt.Println("Cool conditions")
default:
    fmt.Println("Warm conditions")
}
\end{lstlisting}

This form is useful when different conditions must be evaluated
rather than comparing one expression against a set of values.

\subsection{The \texttt{for} Loop}

The \texttt{for} statement is Go's principal iteration
construct. It supports repeated execution of a block of
statements and can represent conventional counter-controlled
loops, condition-controlled loops, and infinite loops.
Go does not have separate \texttt{while} or \texttt{do-while}
statements.

\subsubsection{Traditional Three-Clause Loop}

A three-clause \texttt{for} loop contains an initialisation
statement, a Boolean condition, and a post statement. Its
general structure is:

\begin{lstlisting}
for initialisation; condition; post {
    // Repeated statements
}
\end{lstlisting}

The initialisation statement executes once before the loop
begins. Before each iteration, the condition is evaluated.
If it is true, the loop body executes, followed by the post
statement. The process continues until the condition evaluates
to \texttt{false}.

For example:

\begin{lstlisting}
for i := 0; i < 5; i++ {
    fmt.Println(i)
}
\end{lstlisting}

This loop prints the integers from 0 to 4. The variable
\texttt{i} is initialised to zero, the condition checks
whether it is less than five, and the increment statement
increases its value after each iteration.

\subsubsection{Condition-Only Loop}

A \texttt{for} statement can contain only a condition,
providing functionality similar to a conventional
\texttt{while} loop.

\begin{lstlisting}
count := 5

for count > 0 {
    fmt.Println(count)
    count--
}
\end{lstlisting}

The loop continues while \texttt{count} is greater than zero.
When the condition becomes false, execution proceeds to the
statement following the loop.

\subsubsection{Infinite Loop}

An infinite loop is created by omitting the condition and
the other clauses from the \texttt{for} statement.

\begin{lstlisting}
for {
    // Repeated operations
}
\end{lstlisting}

Such loops are useful in server processes, event-processing
systems, and applications that must continue running until
an explicit termination condition is encountered. They
should include an appropriate mechanism for terminating
execution when required.

\subsubsection{Iteration Using \texttt{range}}

The \texttt{range} clause allows iteration over supported
data structures, including arrays, slices, strings, maps,
and channels. It provides a convenient way to access
elements or values without manually managing an index.

For example, iteration over a slice can be implemented as
follows:

\begin{lstlisting}
numbers := []int{10, 20, 30, 40}

for index, value := range numbers {
    fmt.Println(index, value)
}
\end{lstlisting}

In this example, \texttt{index} represents the element's
index, while \texttt{value} represents the corresponding
element. The loop executes once for each element in the
slice.

If only the element value is required, the index can be
discarded using the blank identifier:

\begin{lstlisting}
for _, value := range numbers {
    fmt.Println(value)
}
\end{lstlisting}

The exact behaviour of \texttt{range} depends on the
underlying data structure. For example, ranging over a map
does not guarantee a particular iteration order, while
ranging over a channel receives values until the channel
is closed.

\subsection{Control Transfer Statements}

Control transfer statements alter the normal sequential
execution of a program. Go provides \texttt{break},
\texttt{continue}, \texttt{goto}, and \texttt{return} for
different forms of control transfer.

\subsubsection{The \texttt{break} Statement}

The \texttt{break} statement terminates the innermost
enclosing \texttt{for}, \texttt{switch}, or \texttt{select}
statement, unless a label identifies another enclosing
construct that can be terminated.

For example:

\begin{lstlisting}
for i := 0; i < 10; i++ {
    if i == 5 {
        break
    }
    fmt.Println(i)
}
\end{lstlisting}

The loop terminates when \texttt{i} reaches five. Therefore,
the values from zero to four are printed.

\subsubsection{The \texttt{continue} Statement}

The \texttt{continue} statement skips the remaining statements
in the current iteration of the innermost enclosing
\texttt{for} loop and proceeds to the next iteration.

\begin{lstlisting}
for i := 1; i <= 5; i++ {
    if i == 3 {
        continue
    }
    fmt.Println(i)
}
\end{lstlisting}

The value three is omitted from the output because the
\texttt{continue} statement skips the print operation for
that iteration. The loop continues with the subsequent
values.

\subsubsection{The \texttt{goto} Statement}

The \texttt{goto} statement transfers control to a labelled
statement within the same function. Its use is restricted
by Go's language specification, including rules that prevent
jumps into the scope of variables that have not yet been
declared at the destination.

\begin{lstlisting}
package main

import "fmt"

func main() {
    i := 0

Loop:
    if i < 3 {
        fmt.Println(i)
        i++
        goto Loop
    }
}
\end{lstlisting}

The label \texttt{Loop} identifies the destination of the
jump. The program repeatedly returns to this point until
the condition becomes false. Although valid in suitable
contexts, \texttt{goto} should generally be used sparingly,
as conventional loops and functions often provide clearer
control structures.

\subsubsection{The \texttt{return} Statement}

The \texttt{return} statement terminates the execution of
the current function and optionally provides values to
the caller. It is frequently used to return computed
results or terminate execution when a condition is met.

\begin{lstlisting}
func divide(a, b float64) (float64, bool) {
    if b == 0 {
        return 0, false
    }

    return a / b, true
}
\end{lstlisting}

The function returns a numerical result and a Boolean
status indicator. When the divisor is zero, it returns
immediately without performing the division. Otherwise,
it returns the quotient and a successful status value.

\subsection{The \texttt{select} Statement}

The \texttt{select} statement coordinates communication
operations on channels, which are Go's built-in mechanism
for communicating values between concurrent activities.
It allows a goroutine to wait for multiple channel
operations and proceed when one of them can be performed.

A \texttt{select} statement has a structure similar to
a \texttt{switch}, but its cases contain channel send or
receive operations.

\begin{lstlisting}
select {
case value := <-channelA:
    fmt.Println("Received:", value)
case channelB <- 10:
    fmt.Println("Value sent")
default:
    fmt.Println("No communication is ready")
}
\end{lstlisting}

When one or more communications are ready, the
\texttt{select} statement chooses one ready case and
executes its associated statements. If multiple cases
are ready, one is selected pseudo-randomly. If no case
is ready and a \texttt{default} clause is present, the
default statements execute immediately. Without a
\texttt{default} clause, the goroutine blocks until a
communication can proceed.

The \texttt{select} statement is particularly important
in concurrent programs that need to coordinate multiple
channels, implement time-outs, or respond to cancellation
signals.

\subsection{Deferred Function Execution}

The \texttt{defer} statement schedules a function call
to execute when the surrounding function returns, either
through an explicit \texttt{return} statement or by
reaching the end of its body. Deferred calls also execute
during normal stack unwinding caused by a panic.

Deferred function calls are executed in last-in, first-out
order. This means that the most recently deferred function
is executed first.

\begin{lstlisting}
func process() {
    defer fmt.Println("First deferred call")
    defer fmt.Println("Second deferred call")

    fmt.Println("Processing")
}
\end{lstlisting}

The output is:

\begin{lstlisting}
Processing
Second deferred call
First deferred call
\end{lstlisting}

This behaviour makes \texttt{defer} useful for releasing
resources, closing files, unlocking mutexes, and performing
cleanup operations. The function value and arguments of a
deferred call are evaluated when the \texttt{defer}
statement executes, although the deferred function itself
runs later.

\subsection{Panic and Recovery}

Go provides \texttt{panic} and \texttt{recover} for handling
abnormal run-time conditions and controlling panic
propagation. A panic interrupts ordinary execution and
initiates the unwinding of the current goroutine's call
stack. During this process, deferred functions are executed.

The built-in \texttt{recover} function can intercept a
panic when called directly from a deferred function in
the same goroutine.

\begin{lstlisting}
func safeOperation() {
    defer func() {
        if r := recover(); r != nil {
            fmt.Println("Recovered from panic:", r)
        }
    }()

    panic("unexpected condition")
}
\end{lstlisting}

In this example, the deferred function invokes
\texttt{recover} and detects the panic value. Execution
then continues according to the program's subsequent
control flow after \texttt{safeOperation} returns.

Panic and recovery should generally be reserved for
exceptional conditions rather than used as replacements
for ordinary error handling. Go programs commonly
represent expected failures using returned \texttt{error}
values.

\subsection{Scope and Execution of Control Structures}

Control flow constructs in Go interact with lexical scope.
Variables declared within an initialisation statement,
a loop body, or a conditional block are subject to the
scope rules of the enclosing construct. This helps limit
the visibility of temporary values and reduces unintended
interactions between different parts of a program.

For example:

\begin{lstlisting}
if value := 10; value > 5 {
    fmt.Println(value)
}

// The variable value is not accessible here.
\end{lstlisting}

The variable \texttt{value} is available within the
\texttt{if} statement, including its associated branches,
but is not accessible after the statement has completed.

The execution of control structures also depends on
their nesting. For example, a \texttt{break} statement
normally terminates the innermost eligible enclosing
loop or selection construct, while a \texttt{return}
statement terminates the entire current function.
Labels can modify the targets of certain control transfer
statements where the language specification permits it.

\subsection{Conclusion}

Control flow in Go provides a structured and expressive
framework for determining the order of program execution.
Conditional statements support decision-making, switch
statements enable multi-way selection, and the versatile
\texttt{for} statement implements different forms of
iteration. Control transfer statements alter the normal
execution path, while \texttt{select} supports channel-based
communication in concurrent programs.

Furthermore, deferred function execution and panic recovery
provide mechanisms for resource management and exceptional
run-time behaviour. An understanding of these constructs,
their scope rules, and their interaction is essential for
developing reliable, efficient, and maintainable Go
applications.